\documentclass{article}
\usepackage{graphicx}
\usepackage{amsmath}
\usepackage{amssymb}
\usepackage{amsfonts}
\usepackage{hyperref}
\usepackage{url}
\usepackage{listings}
\usepackage{xcolor}
\usepackage{float}

\lstset{
    language=Python,
    basicstyle=\ttfamily\small,
    keywordstyle=\color{blue},
    stringstyle=\color{red},
    commentstyle=\color{gray},
    showstringspaces=false,
    frame=single,
    breaklines=true
}



\title{Reconocimiento de Patrones (ML) - T05}
\author{ALEJANDRO ZARATE MACIAS}
\date{2 de Marzo 2026}

\begin{document}

\maketitle

\begin{abstract}
Estos ejercicios tienen como objetivo revisar algunos modelos de aprendizaje no supervisado.
\end{abstract}


% ========================================
% SECCIÓN 1
% ========================================
\section{Problema 1}

\subsection{Enunciado}

Lea \cite{featureSelectionMethods} sobre los métodos de la familia de selección de características. Luego, escriba su criterio personal para elegir uno sobre otro.

\subsection{Análisis}

% TODO

\subsection{Conclusión}

% TODO


% ========================================
% SECCIÓN 2
% ========================================
\section{Problema 2}

\subsection{Enunciado}

Considere la siguiente función:
\[
    f(x, y) = \frac{RMSE(x, y)}{\|x\|_2}, \quad x, y \in \mathbb{R}^n, \quad \|x\|_2 \neq 0.
\]

Escriba el dominio e imagen de $f$. ¿Cuál es un ``buen'' valor/rango para $f$ de modo que podamos decir que $x$ es muy similar a $y$?

\subsection{Análisis}

% TODO

\subsection{Conclusión}

% TODO


% ========================================
% SECCIÓN 3
% ========================================
\section{Problema 3}

\subsection{Enunciado}

Considere el dataset de pronóstico de divorcio \cite{kaggleDivorcePrediction}. A continuación:
\begin{itemize}
    \item [a)] Resuelva el problema de clasificación.
    \item [b)] Identifique las características más relevantes para la clasificación. El objetivo aquí es reducir el número de características empleadas manteniendo una alta precisión de clasificación.
\end{itemize}
Anote todos sus supuestos y discuta sus hallazgos.

\subsection{Metodología}

% TODO

\subsection{Resultados}

% TODO

\subsection{Discusión}

% TODO

\subsection{Conclusión}

% TODO


% ========================================
% SECCIÓN 4
% ========================================
\section{Problema 4}

\subsection{Enunciado}

Repita el problema 3, pero esta vez cree nuevas características. El objetivo es el mismo que antes: reducir el número de características utilizadas manteniendo una alta precisión de clasificación. ¿Fue posible disminuir el número de características usando este enfoque? Anote todos sus supuestos y discuta sus hallazgos.

\subsection{Metodología}

% TODO

\subsection{Resultados}

% TODO

\subsection{Discusión}

% TODO

\subsection{Conclusión}

% TODO


% ========================================
% SECCIÓN 5
% ========================================
\section{Problema 5}

\subsection{Enunciado}

Repita el problema 3 considerando el problema de pronóstico de precio de automóviles \cite{kaggleCarPricePrediction}.

\subsection{Metodología}

% TODO

\subsection{Resultados}

% TODO

\subsection{Discusión}

% TODO

\subsection{Conclusión}

% TODO


% ========================================
% SECCIÓN 6
% ========================================
\section{Problema 6}

\subsection{Enunciado}

Repita el problema 4 considerando el problema de pronóstico de precio de automóviles \cite{kaggleCarPricePrediction}.

\subsection{Metodología}

% TODO

\subsection{Resultados}

% TODO

\subsection{Discusión}

% TODO

\subsection{Conclusión}

% TODO


% ========================================
% SECCIÓN 7
% ========================================
\section{Problema 7}

\subsection{Enunciado}

Utilice los algoritmos K-Means y DBSCAN para realizar compresión de imágenes por colores, como se explica en \cite{imageCompressionKMeans}. Anote todos sus supuestos, haga comparaciones y discuta sus hallazgos.

\subsection{Metodología}

% TODO

\subsection{Resultados}

% TODO

\subsection{Discusión}

% TODO

\subsection{Conclusión}

% TODO


% ========================================
% SECCIÓN 8
% ========================================
\section{Problema 8}

\subsection{Enunciado}

¿Es K-Means demasiado sensible a los valores atípicos debido al operador de media? Averigüemos:
\begin{itemize}
    \item [a)] Construya un dataset con algunos valores atípicos que muestre qué tan sensible es K-Means.
    \item [b)] Codifique un algoritmo de clustering K-Mediana, que es básicamente K-Means pero usa la mediana en lugar de la media.
    \item [c)] Aplique su K-Mediana a sus datos. ¿Agrupa mejor?
\end{itemize}
Anote todos sus supuestos y discuta sus hallazgos.

\subsection{Metodología}

% TODO

\subsection{Resultados}

% TODO

\subsection{Discusión}

% TODO

\subsection{Conclusión}

% TODO


% ========================================
% SECCIÓN 9
% ========================================
\section{Problema 9}

\subsection{Enunciado}

Considere el problema de predicción de género \cite{kaggleWeightHeightOutlier}. Haga lo siguiente:
\begin{enumerate}
    \item Realice un proceso de clustering sobre las características considerando dos clusters. ¿Sus resultados se alinean con las etiquetas conocidas?
    \item Intente identificar todos los valores atípicos.
    \item Seleccione algunos datos aleatorios e invierta su etiqueta. Luego, intente identificar todos los valores atípicos. ¿Su método fue capaz de identificar exitosamente los datos invertidos?
\end{enumerate}
Anote todos sus supuestos, haga comparaciones y discuta sus hallazgos.

\subsection{Metodología}

% TODO

\subsection{Resultados}

% TODO

\subsection{Discusión}

% TODO

\subsection{Conclusión}

% TODO


% ========================================
% REFERENCIAS
% ========================================
\bibliographystyle{plain}
\bibliography{ref}

\end{document}
